\chapter*{Chapitre 6 — Climatisation et traitement d'air}
\addcontentsline{toc}{chapter}{Chapitre 6 — Climatisation et traitement d'air}

% ============================================================
\section*{Énoncés}
% ============================================================

\exercice{Analyse du cycle frigorifique R454B}

Un climatiseur au R454B présente les enthalpies suivantes mesurées au banc d'essai :
$h_1 = \SI{490}{\kilo\joule\per\kilogram}$, $h_2 = \SI{535}{\kilo\joule\per\kilogram}$, $h_3 = \SI{270}{\kilo\joule\per\kilogram}$, $h_4 = h_3$.

Le débit massique mesuré est $\dot{m} = \SI{0{,}08}{\kilogram\per\second}$.

\begin{enumerate}
  \item Calculer les puissances frigorifique, de compression et de condensation.
  \item Vérifier le bilan énergétique.
  \item Calculer l'EER. Comparer avec l'EER de Carnot ($T_{0} = 7$\,°C, $T_{c} = 45$\,°C).
  \item Le fabricant annonce EER = 4,1 dans ces conditions. Analyser l'écart.
\end{enumerate}

\exercice{Charges thermiques d'un open-space et sélection d'un système}

Un open-space de 200\,m² est situé en façade ouest. Données estivales :
\begin{itemize}
  \item Vitrage ouest : 30\,m², $g = 0{,}25$ (vitrage solaire), $F_{ombr} = 0{,}6$, $G_{sol,ouest} = \SI{450}{\watt\per\square\metre}$
  \item 40 personnes, 80\,W/pers
  \item Éclairage : 8\,W/m²
  \item Informatique : 40 postes à 120\,W
  \item Transmission et ventilation : \SI{3000}{\watt}
\end{itemize}

\begin{enumerate}
  \item Calculer chaque type d'apport et la puissance frigorifique totale.
  \item Proposer deux solutions techniques (une à fluide frigorigène, une à eau glacée) en justifiant le choix.
  \item Calculer la consommation électrique mensuelle du système à eau glacée si EER = 4,0 et si le système fonctionne 8\,h/jour, 22\,jours/mois à plein régime.
\end{enumerate}

\exercice{Diagnostic d'un groupe froid sous-performant}

Un groupe froid à eau glacée de 80\,kW nominal présente les mesures suivantes lors de la mise en service :
\begin{itemize}
  \item Puissance électrique mesurée : \SI{28}{\kilo\watt}
  \item Température eau glacée départ : 9\,°C (consigne : 7\,°C)
  \item Température eau glacée retour : 13\,°C
  \item Débit eau glacée : \SI{30}{\cubic\metre\per\hour}, $\rho = \SI{999}{\kilogram\per\cubic\metre}$, $c_p = \SI{4182}{\joule\per\kilogram\per\kelvin}$
  \item Température extérieure : 28\,°C
\end{itemize}

\begin{enumerate}
  \item Calculer la puissance frigorifique produite.
  \item Calculer l'EER mesuré.
  \item Identifier les anomalies et les causes probables.
\end{enumerate}

\newpage
% ============================================================
\section*{Corrigés}
% ============================================================

\subsection*{Corrigé — Exercice 1}

\textit{Question 1 — Puissances :}
\[
  q_0 = 490 - 270 = \SI{220}{\kilo\joule\per\kilogram}
\]
\[
  w_c = 535 - 490 = \SI{45}{\kilo\joule\per\kilogram}
\]
\[
  q_c = 535 - 270 = \SI{265}{\kilo\joule\per\kilogram}
\]
\[
  \dot{Q}_0 = 0{,}08 \times 220 = \boxed{\SI{17{,}6}{\kilo\watt}}, \quad P_{élec} = 0{,}08 \times 45 = \boxed{\SI{3{,}6}{\kilo\watt}}, \quad \dot{Q}_c = 0{,}08 \times 265 = \boxed{\SI{21{,}2}{\kilo\watt}}
\]

\textit{Question 2 — Bilan :} $q_0 + w_c = 220 + 45 = 265 = q_c$ \checkmark

\textit{Question 3 — EER :}
\[
  EER_{mes} = \frac{17{,}6}{3{,}6} = \boxed{4{,}89}
\]
\[
  EER_{Carnot} = \frac{T_0}{T_c - T_0} = \frac{280}{318 - 280} = \frac{280}{38} = \boxed{7{,}37}
\]
\[
  \eta_{ex} = \frac{4{,}89}{7{,}37} = 66{,}4\,\%
\]

\textit{Question 4 — Écart EER :} L'EER mesuré (4,89) est supérieur à l'EER annoncé (4,1). L'écart s'explique par : des conditions d'essai différentes des conditions normalisées EN 14825, l'instrumentation de mesure du débit massique, et les incertitudes de mesure des enthalpies. L'EER annoncé intègre les pertes auxiliaires (ventilateurs, régulation) que le calcul enthalpique n'inclut pas.

\subsection*{Corrigé — Exercice 2}

\textit{Question 1 — Apports :}
\[
  Q_{sol} = 30 \times 450 \times 0{,}25 \times 0{,}6 = \SI{2025}{\watt}
\]
\[
  Q_{pers} = 40 \times 80 = \SI{3200}{\watt}
\]
\[
  Q_{écl} = 8 \times 200 = \SI{1600}{\watt}
\]
\[
  Q_{info} = 40 \times 120 = \SI{4800}{\watt}
\]
\[
  \dot{Q}_{froid} = 2025 + 3200 + 1600 + 4800 + 3000 = \boxed{\SI{14\,625}{\watt} \approx \SI{14{,}6}{\kilo\watt}}
\]

\textit{Question 2 — Solutions :}
\begin{itemize}
  \item \textbf{VRV/VRF (fluide frigorigène)} : adapté si bâtiment multi-zones, flexibilité élevée, pas besoin de réseau eau glacée. Recommandé pour réhabilitation ou bâtiment de taille moyenne.
  \item \textbf{Eau glacée + ventilo-convecteurs 4 tubes} : adapté si bâtiment neuf de grande taille avec besoins chaud et froid simultanés, meilleur EER, maintenance centralisée. Recommandé pour immeuble tertiaire neuf.
\end{itemize}

\textit{Question 3 — Consommation électrique mensuelle :}
\[
  P_{élec} = \frac{\dot{Q}_{froid}}{EER} = \frac{14{,}6}{4{,}0} = \SI{3{,}65}{\kilo\watt}
\]
\[
  E_{mensuelle} = 3{,}65 \times 8 \times 22 = \boxed{\SI{642}{\kilo\watt\heure}}
\]

\subsection*{Corrigé — Exercice 3}

\textit{Question 1 — Puissance frigorifique :}
\[
  \dot{m} = \frac{30 \times 999}{3600} = \SI{8{,}33}{\kilogram\per\second}
\]
\[
  \dot{Q}_0 = 8{,}33 \times 4182 \times (13 - 9) = \boxed{\SI{139}{\kilo\watt}}
\]

\begin{attention}
La puissance produite (139\,kW) est largement supérieure à la puissance nominale (80\,kW), ce qui constitue la première anomalie à analyser.
\end{attention}

\textit{Question 2 — EER mesuré :}
\[
  EER = \frac{139}{28} = \boxed{4{,}96}
\]

\textit{Question 3 — Anomalies :}
\begin{itemize}
  \item La puissance produite (139\,kW) est \textbf{supérieure} à la puissance nominale (80\,kW) → le groupe est en surcharge ou les conditions sont plus favorables que prévues (température extérieure plus basse, ou nominal mesuré à des conditions défavorables).
  \item La température de départ (9\,°C) est \textbf{au-dessus} de la consigne (7\,°C) → vérifier la régulation, le sous-refroidissement insuffisant, ou un débit trop élevé.
  \item \textbf{Pistes d'action :} recalibrer la vanne de régulation du débit ; vérifier la charge en fluide frigorigène ; vérifier la consigne du thermostat de départ.
\end{itemize}
